\section{Static models}
I use a stylised 3-factor production model throughout this analysis.
$Y$, $A$, $K$, and $L$ are each firm-year observation's structural GVA, TFP, labour, and capital as usual.
I use lower-case variables to indicate logged terms in reduced form
and bolded variables to indicate stacked vectors.
\par
Log TFP, denoted by $z_{i,t}$, is the reduced-form residual after controlling for the two classical inputs,
which are assumed to be exogenous and exhibit unit factor substituability in this static model.
Under this specification, the residual is equivalent to the sum of a constant across all observations,
firm and time fixed effects, and most crucially, the weighted sum of all peer log TFP $\sum_{j\neq i} w_{ij}z_{j,t}$,
with a structural $\delta_0$ parameter to estimate.

\begin{align}
    Y_{i,t} &= A_{i,t} \cdot
                K_{i,t}^{\beta_1} \cdot
                L_{i,t}^{\beta_2}
                % \cdot \left[  \prod_{j\neq i}Z_{j,t}^{w_{ij}} \right]^\gamma
    \\
    \label{eq:static-main-gva}
    y_{i,t} &= \alpha_i + \gamma_t + \beta_0
                + \beta_1 k_{i,t}
                + \beta_2 l_{i,t}
                + \delta_0 \sum_{j\neq i} w_{ij}z_{j,t}
                + \varepsilon_{i,t}
    \\ \notag
    z_{i,t} &= \ln(A_{it}) = y_{i,t} - \alpha l_{i,t} - \beta k_{i,t} 
    \\
    \label{eq:static-main-reduced}
    z_{i,t} &= \alpha_i + \gamma_t + \beta_0 + \delta_0 \sum_{j\neq i} w_{ij}z_{j,t} + \varepsilon_{i,t}
\end{align}

Theory suggests no restrictions on $\beta_1$, $\beta_2$, and $\delta_0$ collectively should be imposed,
as $\delta_0$ is an externality. \todo{Lucas discusses this}
\par
The reduced-form equation \eqref{eq:static-main-gva} is linear in the weights $w_{ij}$
applied to peer log TFP and its coefficient.
This permits me to estimate $\delta_0$ through either
the full GVA specification in one stage, using the log of \itx{gva1} as the dependent variable,
or in two stages, calculating $z_{i,t}$ as the 2-factor residual
and then regressing $z_{i,t}$ however desired in reduced form.
This is convenient for a number of the static models below.

\subsection{2-factor TFP estimation}

For convenience, I initially conduct the 2-factor estimation
deriviing the residual $z_{ijt}$ to decompose further in a second stage.
This step is additionally helpful for validating the relationships about
the derived aggregates from my data.

\begin{align}
    \label{eqn:tfp2factor}
    y_{i,t} &= \alpha_i + \gamma_t + \beta_0
                + \beta_1 k_{i,t}
                + \beta_2 l_{i,t}
                + u_{i,t}
    \\ \notag
    u_{i,t} &= \delta_0 \sum_{j\neq i} w_{ij}z_{j,t} + \varepsilon_{i,t}
    \\
    z_{i,t} &= \alpha_i + \gamma_t + \beta_0 + u_{i,t}
    \\
    \hat z_{i,t} &= y_{i,t} - \hat\beta_1k_{i,t} - \hat\beta_2l_{i,t}
\end{align}
%
I conduct the estimation using the log-transformed variables.
I must include firm and time fixed effects and a constant
if equation (\ref{eq:static-main-reduced}) is the true specification,
otherwise the static panel OLS will inaccurately attribute firm and time variation
to my classical factor inputs \todo{explain more clearly}.
As outlined above, this analysis assumes factor input exogeneity.

\wraptable
    {Baseline TFP calculation from 2-factor production}
    {../../model/output/results_0_tfp}
    {
        All parameters reported with standard errors underneath.
        Significant estimates against an $H_0: \beta=0$ are indicated
        with $p<0.1$ (*), $p<0.05$ (**), and $p<0.01$ (***).
    }

The parameter coefficients give a labour income share exactly in line with ONS analysis of UK industry-level data
over the same time period, particularly \itx{corporate gross labour share} which is constructed
from a similar sample to the coverage in Fame\footnote{\cite{OfficeforNationalStatistics2024}}.
\par
Surprisingly, running the regression with \itx{fixed\_total} as the capital measure instead of \itx{total\_assets}
almost entirely eliminates the capital share of income, or explanatory power of capital on variation in GVA.
I leave it for future analysis why this dataset yields this phenomenon and
for the rest of the analysis, I use \itx{total\_assets} as the primary capital measure.
\par
As a demonstration, I include specifications which omit fixed effects.
The capital share of income jumps on a specification with only time fixed effects,
which is to be expected \todo{explain}.

TFP is the residual on GVA after extracting the proportion that can be accounted by capital and labour inputs.
Equivalently, it is the sum of fixed effects in this specification with the error term,
which means we should expect it to be exogenous to capital and labour inputs,
but correlated to both firm effects and time trends.
By extracting this regression residual, I have a TFP measure for all of my 1,081,520 panel observations.

% ------

\subsection{Static panel identification}
\label{ss-static-identification}

Models in this section are consistent only if
the regressors are plausibly exogenous from $\varepsilon_{i,t}$.
Capital and labour are assumed to be exogenous in this static model;
we can broadly interpret this as inflexible factor input decisions.
\begin{align}
    \notag
    x_{i,t} &=
    \begin{pmatrix}
        k_{i,t} &   l_{i,t} &   \sum_{j\neq i} w_{ij,t}z_{j,t}
    \end{pmatrix}
    \\ \notag
    \mathbb E[\varepsilon_{i,t} \vert k_{i,t}] &= 0,\ 
    \mathbb E[\varepsilon_{i,t} \vert l_{i,t}] = 0
\end{align}

This leaves only the argument for the exogeneity of $\sum_{j\neq i} w_{ij,t}z_{j,t} \in x_{i,t}$,
the peer effect regressor.
Stacking $w_{ij}$ target-peer weights to an $n\times n$ matrix of firm-firm interactions $\textbf W$
yields a spatial econometric reflection problem,
comparable to the MA inversion procedure in autoregressive time series analysis.

\begin{align}
    \eqreref{eq:static-main-reduced}
    z_{j,t} &= \alpha_j + \gamma_t + \beta_0 + \delta_0 \sum_{k\neq j} w_{jk}z_{k,t} + \varepsilon_{j,t}
    \\ \notag
    \mathbf z_t &= \boldsymbol\alpha + \gamma_t \boldsymbol\iota + \beta_0\boldsymbol\iota + \delta_0 \mathbf W \mathbf z_t + \boldsymbol\varepsilon_t
    \\ \notag
                &= (\mathbb I-\delta_0 \mathbf W)^{-1}(\boldsymbol\alpha + \gamma_t \boldsymbol\iota + \beta_0\boldsymbol\iota + \boldsymbol\varepsilon_t)
    \\
    \label{eq:static-z-ma-inft}
                &= \sum_{k=0}^\infty \delta_0^k \mathbf W^k \cdot (\boldsymbol\alpha + \gamma_t \boldsymbol\iota + \beta_0\boldsymbol\iota + \boldsymbol\varepsilon_t)
    \\ \notag
    \mathbb E[
        \varepsilon_i
        \vert
        \boldsymbol W^k
        \begin{pmatrix}
            \boldsymbol\alpha   &   \gamma_t \boldsymbol\iota   &  \beta_0\boldsymbol\iota
        \end{pmatrix}
    ] &=0,\  \forall k
    \\
    \label{eq:static-ass-ma-exogeneity}
    \mathbb E[
        \varepsilon_{i,t}
        \vert
        \boldsymbol W^k
        \boldsymbol \varepsilon_t
    ] &=0,\  
    \forall k
    \\
    \label{eq:static-ass-error-orthogonality}
    \implies \mathbb E[\epsilon_{i,t}\epsilon_{j,t}] &= 0,\ \forall k
\end{align}
%
The reflection problem shows that nearly all components of peer effect regressor $\sum_{j\neq i}w_{ij}z_{j,t}$,
namely the time-invariant weighting $w_{ij}$, fixed effects, and the constant are unconfounded.
However, the weighted peer TFP term contains its own idiosyncratic error
which is weighted by the sum of $\mathbf W$ to every power by the spatial autocorrelation inversion.
A strong exogeneity assumption is required,
$\mathbb E[\varepsilon_{i,t}\vert\boldsymbol W^k\boldsymbol \varepsilon_t] =0,\ \forall k]$,
which in itself implies the already-strong assumption
of orthogonality between idiosyncratic errors hitting peer firms \eqref{eq:static-ass-error-orthogonality}
This rules out any structural, cross-sectional correlation between error terms:
any productivity shock $\epsilon_{j,t}$ occurring in peer firm $j$,
if it is correlated with a productivity shock $\epsilon_{i,t}$
occurring in target firm $i$ contemperaneously must be attributable to the peer effect coefficient $\delta_0$.
\par
Assumption \eqref{eq:static-ass-error-orthogonality} is not sufficient for $\mathbb E[\epsilon_{i,t}|z_{j,t}]=0$,
as I discuss in section \ref{ss:model-dd}.
Even if were, it is quite a strong assumption. It precludes any contemperaneous, correlated independent TFP shocks,
which might occur on both the global and local level. \todo{...}


% --- %

\subsection{Linear-in-means grouped peers}
\label{ss:model-lmm}

Assuming exogeneity assumptions \eqref{eq:static-ass-ma-exogeneity} and \eqref{eq:static-ass-error-orthogonality} hold,
the following reduced-form equation can be estimated with a simple panel OLS.

\begin{align}
    \eqreref{eq:static-main-reduced}
    z_{i,t} &= \alpha_i + \gamma_t + \beta_0 + \delta_0 \sum_{j} w_{ij}z_{j,t} + \varepsilon_{i,t}
    \\
    w_{ij}\in \mathbf W &=
        \begin{cases}
            0,              & \text{if}         \ i=j           \\
            \frac1{n_g-1},  & \text{else if}    \ (i,j)\in g    \\
            0,              & \text{otherwise}
        \end{cases}
    \\ \notag
    &\sim
    \begin{pmatrix}
        0           & \frac{1}{3} & \frac{1}{3} & \frac{1}{3} & 0           & 0           & 0           \\
        \frac{1}{3} & 0           & \frac{1}{3} & \frac{1}{3} & 0           & 0           & 0           \\
        \frac{1}{3} & \frac{1}{3} & 0           & \frac{1}{3} & 0           & 0           & 0           \\
        \frac{1}{3} & \frac{1}{3} & \frac{1}{3} & 0           & 0           & 0           & 0           \\
        0           & 0           & 0           & 0           & 0           & \frac{1}{2} & \frac{1}{2} \\
        0           & 0           & 0           & 0           & \frac{1}{2} & 0           & \frac{1}{2} \\
        0           & 0           & 0           & 0           & \frac{1}{2} & \frac{1}{2} & 0           
    \end{pmatrix}
\end{align}
%
The initial static model applies simple time-invariant weights $w_{ij}$
of the averages of spatial peers.
As has been indicated by the sum of peers $j\neq i$,
I adopt a 'leave-one-out' approach, where the resulting $n\times n$ matrix
is zero on the diagonal and the reciprocal weight applied across peers
is therefore the number of group members $n_g-1$, which excludes the target firm $i$.
Intuitively, this is the unweighted average of peer TFP within a spatial ring.
The specification is known as the 'Linear-in-Means (LMM)' Model
and is common to the peer effects literature.
I illustrate an example of a weighted matrix $W$ following this rule for 7 firms,
categorised by a group of 4 and a group of 3.
\par
I discussed group assignment in section \ref{ss:data-distance}.
Firms have been labelled using their location features to symmetric, transitive,
non-overlapping groups by \itx{ttwa} (largest) to \itx{lat\_lon5} (smallest),
with the latter meaning two firm entries shared a latitude-longitude identifier
to 5 decimal places.
Reviewing the summary statistics, I opted for full postcodes, \itx{pc8}
as the group identifier for Model (1).
This was mostly driven by full data availability, and the consistency of postcode definitions
with the larger group $pc4$. Deriving a subgroup by latitude-longitude match
within the postcode seemed fruitless, as in many cases it would identify the exact same firms,
a move from the average number of peers from 3.1 to 2.9,
and was a similar number of groups, with non-coheisve definitions.
\par
\begin{table}[htbp]
    \centering
    \caption{\todo{Placeholder} Summary statistics across spatial groups}
    \label{tab:spatial_group_summary}
    \begin{tabular}{lcccc}
        \toprule
        Property            & Number of groups & Unclassified & Unclassified (\%) & Avg. peers per group \\
        \midrule
        \itx{ttwa}          & 228    & 567 & 0.4\% & 668.3 \\
        \itx{pc4}           & 2,705  & 0   & 0.0\% & 56.3  \\
        \itx{pc8}           & 49,963 & 0   & 0.0\% & 3.1   \\
        \itx{lat\_lon3}     & 48,540 & 507 & 0.3\% & 3.1   \\
        \itx{lat\_lon4}     & 53,193 & 507 & 0.3\% & 2.9   \\
        \itx{lat\_lon5}     & 53,246 & 507 & 0.3\% & 2.9   \\
        \bottomrule
    \end{tabular}
\end{table}
\noindent
The results below show model (1), where peers and targets share a postcode.
\wraptable
    {Static panel linear-in-means using spatial-group TFP averages}
    {../../model/output/results_1a_lmm}
    {}

\newpage
\par
As an extension, I regress the multiple, non-overlapping spatial groups at once.
Multiple spatial autoregressors do not change the reflection problem:
the moving average inversion's coefficient merely becomes the weighted sum
of 3 independent matrices $\delta_g \mathbf W_g$, for a generic set of weights $\mathbf W_g$.
\begin{align}
    \label{eq:lmm-structural1}
    z_{i,t} &= \alpha_i + \gamma_t + \beta_0
                + \sum_{g=1}^3 \delta_g \sum_{j\neq i} w_{ij,g}z_{j,t}
                + \varepsilon_{it}
    \\
    \mathbf z_t &= \sum_{k=0}^\infty (\delta_1 \mathbf W_1 + \delta_2 \mathbf W_2 + \delta_3 \mathbf W_3)^k \cdot
        (\boldsymbol\alpha + \gamma_t \boldsymbol\iota + \beta_0\boldsymbol\iota + \boldsymbol\varepsilon_t)
\end{align}

The weighting rules for these multiple, non-overlapping groups is to zero \todo{...}
These matrices are orthogonal, justifying the result found above that the coefficient on $\delta_0$
is identical in the 1-regressor model $delta_0$ to the 3-regressor model for $\delta_1$.
\begin{align}
  w_{ij}\in \mathbf{W}_g &=
    \begin{cases}
      0,             & \text{if } i=j \\
      0,             & \text{if } w_{ij,k} \neq 0, \exists k<g \\
      \frac{1}{n_g-1}, & \text{else if } (i,j)\in g \\
      0,             & \text{otherwise}
    \end{cases}
  \\ \notag
  &\sim
    \vcenter{\hbox{\scriptsize$\begin{pmatrix}
            0 & \frac{1}{2} & \frac{1}{2} & 0 & 0 & 0 & 0 & 0 \\
            \frac{1}{2} & 0 & \frac{1}{2} & 0 & 0 & 0 & 0 & 0 \\
            \frac{1}{2} & \frac{1}{2} & 0 & 0 & 0 & 0 & 0 & 0 \\
            0 & 0 & 0 & 0 & 0 & 0 & 0 & 0 \\
            0 & 0 & 0 & 0 & 0 & 0 & 0 & 0 \\
            0 & 0 & 0 & 0 & 0 & 0 & 0 & 0 \\
            0 & 0 & 0 & 0 & 0 & 0 & 0 & 1 \\
            0 & 0 & 0 & 0 & 0 & 0 & 1 & 0
        \end{pmatrix}$}}
    \vcenter{\hbox{\scriptsize$\begin{pmatrix}
            0 & 0 & 0 & \frac{1}{4} & \frac{1}{4} & 0 & 0 & 0 \\
            0 & 0 & 0 & \frac{1}{4} & \frac{1}{4} & 0 & 0 & 0 \\
            0 & 0 & 0 & \frac{1}{4} & \frac{1}{4} & 0 & 0 & 0 \\
            \frac{1}{4} & \frac{1}{4} & \frac{1}{4} & 0 & \frac{1}{4} & 0 & 0 & 0 \\
            \frac{1}{4} & \frac{1}{4} & \frac{1}{4} & \frac{1}{4} & 0 & 0 & 0 & 0 \\
            0 & 0 & 0 & 0 & 0 & 0 & \frac{1}{2} & \frac{1}{2} \\
            0 & 0 & 0 & 0 & 0 & \frac{1}{2} & 0 & 0 \\
            0 & 0 & 0 & 0 & 0 & \frac{1}{2} & 0 & 0
        \end{pmatrix}$}}
    \vcenter{\hbox{\scriptsize$\begin{pmatrix}
            0 & 0 & 0 & 0 & 0 & \frac{1}{8} & \frac{1}{8} & \frac{1}{8} \\
            0 & 0 & 0 & 0 & 0 & \frac{1}{8} & \frac{1}{8} & \frac{1}{8} \\
            0 & 0 & 0 & 0 & 0 & \frac{1}{8} & \frac{1}{8} & \frac{1}{8} \\
            0 & 0 & 0 & 0 & 0 & \frac{1}{8} & \frac{1}{8} & \frac{1}{8} \\
            0 & 0 & 0 & 0 & 0 & \frac{1}{8} & \frac{1}{8} & \frac{1}{8} \\
            \frac{1}{8} & \frac{1}{8} & \frac{1}{8} & \frac{1}{8} & \frac{1}{8} & 0 & 0 & 0 \\
            \frac{1}{8} & \frac{1}{8} & \frac{1}{8} & \frac{1}{8} & \frac{1}{8} & 0 & 0 & 0 \\
            \frac{1}{8} & \frac{1}{8} & \frac{1}{8} & \frac{1}{8} & \frac{1}{8} & 0 & 0 & 0
        \end{pmatrix}$}}
\end{align}

The parameter results suggest a strong correlation between the TFP of a target firm and its peers.
Peers firms in the local postcode have an elasticity of 0.29,
significant at the 1\% level, on the target firm's TFP.
However, since I failed to show exogeneity of $z_{j,t}$ within the regressor,
we cannot describe this as causal. Consider omitted variable bias from a local demand shock \todo{...}

% ---

\newpage
\subsection{Static distance decay}
\label{ss:model-dd}

To improve on the LMM specification, I make a modification where
every peer firm is weighted by its influence on the target by the reciprocal of the 
continuous distance metric calculated in \ref{ss:data-distance}
$w_{ijt}=f(d)$. This is stylistic of a declining peer influence further away,
representing a \todo{lower frequency of social collisions, barriers created via transport and time costs,
and an increased difficulty of observation of management practices}.

\begin{align}
    \label{eq:dd-structural1}
    z_{it} &= \alpha_i + \gamma_t + \beta_0
            + \beta_1 \sum_{j\neq i} w_{ijt}z_{jt} + \varepsilon_{it}
    \\
    \mathbf z_{t} &= (\alpha_i + \gamma_t + \beta_0)\boldsymbol\iota
            + \beta_1 \sum_{j\neq i} \Omega_t \mathbf z_{t} + \boldsymbol\varepsilon_{t}
    \\
    w_{ijt}\in \mathbf \Omega_t &=
        \begin{cases}
            \frac1{d_{ij}},         &\text{(1) gravity}                                      \\
            \frac1{d_{ij}^2}        &\text{(2) N.E.G.} \\
            \exp{-d_{ij}\cdot k}    &\text{(3) exponential decline}
        \end{cases}
\end{align}

This represents a likely more accurate declining effect of productivity spillover,
rather than the special-case group mean where every firm in an enclosed area,
for example, the target's postcode, has an equal influence on the target's TFP.
In addition, it allows for better identification following \cite{Bramoulle2009}.
\todo{explain Bramoulle2009, patterns of }$\Omega_t$
\\*
In essence, the LMM specification only contains 3 non-overlapping groups
from its location tags. The distance-decayed specification
contains as many non-overlapping groups as firms.

\wraptable
    {Static panel using continuous distance-decay weighted TFP average of peers}
    {../../model/output/results_2_dd}
    {}

This distance-decay specification revises down the estimate from peers to a an elasticity of 0.168,
lower than the results given by the LMM specification. This is likely due to \todo{...}
\\
In addition, I introduce controls \todo{...}